\section{Proveedores de IA generativa y prompts utilizados}
\label{sec:prompts}

La tarea LEC06 exige el uso de \textbf{al menos dos proveedores de IA
generativa}, indicar los \textit{prompts} optimizados y justificar por
qué los resultados elegidos son los mejores. El equipo empleó los dos
siguientes, ambos en sus interfaces web:

\begin{itemize}
  \item \textbf{Google Gemini} (versión web, modelo Gemini 2.5).
        Priorizado en tareas que requerían \emph{volumen y cobertura}:
        listado amplio de \textit{stakeholders} y guía operativa de
        entrevistas con representantes de usuarios.
  \item \textbf{Anthropic Claude} (versión web, familia Claude 4.X).
        Priorizado en tareas que requerían \emph{rigor estructural,
        consistencia de formato y sensibilidad ética}: funcionalidad
        para usuarios con discapacidad, cumplimiento de módulos UX/UI
        y casos de uso en formato Cockburn extendido.
\end{itemize}

Las respuestas crudas se versionaron en
\texttt{prompts/respuestas/} para auditoría: cada salida puede
contrastarse contra el prompt que la originó y contra la versión
integrada en el informe.

\subsection{Estrategia general}

Toda conversación con cada IA comenzó con un \textbf{contexto base}
idéntico (\texttt{prompts/00\_contexto\_base.md}) que fija dominio,
restricciones de estilo, idioma, lenguaje centrado en la persona y el
principio de \emph{automatización obligatoria}: cualquier
funcionalidad o requerimiento propuesto debe ser verificable por test,
telemetría o validación de UI.

Para cada uno de los puntos exigidos por la tarea se eligió el
proveedor con mayor probabilidad de producir el mejor resultado, y los
resultados se compilaron manualmente en este informe sin copiar
respuestas de IA en bruto. Las divergencias entre la respuesta del
proveedor y el contenido final del informe se documentan en cada
subsección.

\subsection{Prompt 1 --- Stakeholders (Gemini)}
\paragraph{Archivos} prompt: \texttt{prompts/01\_stakeholders\_gemini.md};
respuesta: \texttt{prompts/respuestas/gemini\_01\_stakeholders.md}.

\paragraph{Por qué Gemini} Su tendencia a producir listas amplias y a
incorporar actores institucionales y regulatorios sin que se le
recuerde explícitamente facilita maximizar cobertura desde el primer
borrador.

\paragraph{Criterio de selección y resultado} La respuesta de Gemini
cubrió los cuatro perfiles con discapacidad como \textit{stakeholders}
primarios separados, identificó la Ley
N\textsuperscript{\textordmasculine}~29733 como marco regulatorio peruano
relevante, y nombró concretamente la Oficina de Empleabilidad UNI y
los proveedores cloud (Render, Fly.io, Vercel). Estos descriptores
específicos se integraron en la sección~\ref{sec:stakeholders}.

\paragraph{Divergencia documentada} Gemini propuso 14
\textit{stakeholders} y omitió cuatro roles que el equipo consideró
relevantes: \textit{peer} entrevistador como rol activo, reclutador
como consultor del proyecto, coach/mentor y administrador del sistema
diferenciado del equipo de desarrollo. La sección 01 los agrega
explícitamente y mantiene su numeración \sk{1}--\sk{17} como
autoritativa, con tabla de equivalencia hacia la numeración de Gemini.

\paragraph{Incidencia detectada} En la ejecución del 29-04-2026, la
salida del prompt 02 fue idéntica a la del prompt 01 en la sesión web
de Gemini (probable reutilización de contexto en la pestaña). El
equipo decidió no reejecutar el prompt 01 y delegar el contenido
sustantivo del punto 2 (funcionalidad para discapacidad) a Claude, que
es donde realmente se requería el rigor de WCAG.

\subsection{Prompt 2 --- Funcionalidad para discapacidad (Claude)}
\paragraph{Archivos} prompt: \texttt{prompts/02\_discapacidad\_claude.md};
respuesta: \texttt{prompts/respuestas/claude\_02\_discapacidad.md}.

\paragraph{Por qué Claude} La sección requiere lenguaje centrado en la
persona, citas correctas a WCAG 2.2 nivel AA y atención a casos donde
una métrica multimodal puede penalizar injustamente al usuario.

\paragraph{Criterio de selección y resultado} La respuesta de Claude
cumplió los cuatro criterios definidos a priori: (i) propuso métricas
sustitutas para cada caso de penalización injusta (postura
$\to$ presencia, contacto visual $\to$ orientación de voz hacia el
micrófono, fluidez verbal $\to$ coherencia y densidad léxica, postura
$\to$ estabilidad del encuadre facial); (ii) citó criterios WCAG
concretos por número (2.1.1, 1.4.3, 1.4.4, 1.2.4, 1.4.13, 1.4.12,
2.5.8, 2.2.1, 2.1.4, 2.5.7, 2.5.4, 2.3.3, 3.2.5, 4.1.2); (iii) usó
lenguaje centrado en la persona en todo el texto; (iv) cerró con un
párrafo de verificación que enumera herramientas (\texttt{axe-core},
Lighthouse, Playwright, panel de telemetría) y cuatro hitos por perfil.
La respuesta se integró íntegramente en la sección~\ref{sec:discapacidad}
con ajustes mínimos de formato LaTeX.

\subsection{Prompt 3 --- Requerimientos funcionales y no funcionales}
\paragraph{Estado} Pendiente de ejecutar formalmente con un proveedor.
El equipo elaboró la sección~\ref{sec:requerimientos} a partir del
\textit{backlog} consensuado en las reuniones de diseño, validado
contra los \textit{stakeholders} de la sección~\ref{sec:stakeholders}
y contra los módulos UX/UI de la sección~\ref{sec:modulos}. El prompt
preparado para Gemini (\texttt{prompts/03\_requerimientos\_gemini.md})
está listo para ejecución en una iteración posterior si el profesor
lo requiere; los criterios de selección ya están definidos en el
prompt.

\subsection{Prompt 4 --- Cumplimiento de módulos UX/UI (Claude)}
\paragraph{Archivos} prompt: \texttt{prompts/04\_modulos\_claude.md};
respuesta: \texttt{prompts/respuestas/claude\_04\_modulos.md}.

\paragraph{Por qué Claude} La sección requiere coherencia
argumentativa entre lo que cada módulo es, cómo lo cumple el
simulador y qué requerimientos satisface.

\paragraph{Criterio de selección y resultado} Se seleccionó la
respuesta de Claude por: (i) precisión mecanística en el mapeo del
aura (color $=$ valencia, intensidad $=$ fluidez, ritmo de pulsación
$=$ cadencia, asimetría espacial $=$ equilibrio verbal/no verbal);
(ii) descripción del \textit{pipeline} de tres roles del LLM con
contexto y memoria acotada propios; (iii) justificación de la
exclusión de hápticas con cuatro argumentos diferenciados (dominio,
hardware, inmersión, recursos), no agrupados en un único párrafo
difuso; (iv) cuatro capas observables de ayuda contextual, no la
clásica enumeración tooltip/asistente. Estos elementos se integraron
en la sección~\ref{sec:modulos}.

\paragraph{Divergencia documentada} Claude usó una numeración de
RF/RNF distinta de la canónica del informe (RF-21..RF-60). Las
referencias en la sección~\ref{sec:modulos} fueron remapeadas a la
numeración \rf{01}--\rf{24} de la sección~\ref{sec:requerimientos}.

\subsection{Prompt 5 --- Casos de uso (Claude)}
\paragraph{Archivos} prompt: \texttt{prompts/05\_casos\_uso\_claude.md};
respuesta: \texttt{prompts/respuestas/claude\_05\_casos\_uso.md}.

\paragraph{Por qué Claude} Para ocho casos de uso largos, la
consistencia entre todos importa más que la cobertura adicional.
Claude respeta plantillas extensas sin fugarse del esquema y, cuando
detecta que falta un RF, lo señala como nota separada al final en
lugar de inventarlo en medio.

\paragraph{Criterio de selección y resultado} Se seleccionó la
respuesta de Claude porque: (i) produjo exactamente los ocho casos de
uso pedidos con la numeración solicitada; (ii) cada CU incluyó al
menos un flujo alternativo documentado; (iii) cerró con dos notas
explícitas sobre cobertura: tres RNF transversales que no encajan en
un CU específico (compatibilidad de navegador, cobertura de pruebas,
consumo de memoria) y dos posibles RF faltantes (matchmaking previo a
peer mock, persistencia de preferencias de accesibilidad entre
dispositivos). Estas observaciones se incorporaron como notas finales
en la sección~\ref{sec:casos}. El cuerpo de los CU del informe usa la
numeración \rf{}/\rnf{} canónica.

\paragraph{Divergencia documentada} Como en el prompt 4, la respuesta
de Claude empleó una numeración propia de RF/RNF (RF-46, RF-60, etc.)
que sugería un \textit{backlog} ampliado. El informe mantiene la
numeración \rf{01}--\rf{24} de la sección~\ref{sec:requerimientos}
como autoritativa; las observaciones de Claude sobre RF faltantes
quedan como insumo para una iteración futura del \textit{backlog}.

\subsection{Prompt 6 --- Guía de entrevistas con usuarios finales (Gemini)}
\paragraph{Archivos} prompt:
\texttt{prompts/06\_entrevistas\_gemini.md}; respuesta:
\texttt{prompts/respuestas/gemini\_06\_entrevistas.md}.

\paragraph{Por qué Gemini} La generación de preguntas variadas para
siete perfiles distintos en una sola pasada se beneficia de la
amplitud de Gemini.

\paragraph{Criterio de selección y resultado} La respuesta cumplió los
criterios: (i) preguntas abiertas en presente sin sesgar la respuesta;
(ii) tiempo total acotado a 45--60 minutos; (iii) bloques específicos
por perfil con discapacidad, no genéricos; (iv) plantilla de
consentimiento informado lista para imprimir. Se integró
íntegramente en la sección~\ref{sec:entrevistas} como guía operativa
para que el equipo ejecute las entrevistas reales.

\subsection{Resumen del flujo de trabajo}

\begin{enumerate}
  \item Pegar el contexto base en una conversación nueva en cada IA.
  \item Esperar la confirmación ``Listo'' antes de enviar el prompt
        específico.
  \item Pegar el prompt del punto correspondiente.
  \item Versionar la respuesta cruda en
        \texttt{prompts/respuestas/<proveedor>\_<prompt>\_<topico>.md}.
  \item Aplicar los criterios de selección descritos en cada
        subsección al resultado.
  \item Si el resultado falla algún criterio, repetir con prompt
        ajustado y conservar las dos versiones.
  \item Compilar manualmente los mejores fragmentos en este informe
        LaTeX y documentar las divergencias.
\end{enumerate}

\subsection*{Apéndice A --- Restricciones del contexto base}
\begin{enumerate}
  \item Español neutro académico, sin emojis ni coloquialismos.
  \item Listas exhaustivas y mutuamente excluyentes.
  \item No inventar datos numéricos; marcar como
        \texttt{(referencial)}.
  \item Toda funcionalidad propuesta debe ser automatizable.
  \item Lenguaje centrado en la persona y referencia a WCAG 2.2 AA.
  \item Salida lista para pegar en LaTeX.
\end{enumerate}
